\subsection{Marco teórico}

El desarrollo de un middleware de orquestación para apoyar procesos de anotación automática en ELAN requiere comprender varios conceptos que se relacionan entre sí. No se trata únicamente de conectar una herramienta de escritorio con un modelo de inteligencia artificial, sino de entender cómo se representan las anotaciones, cómo se procesan los videos, cómo se comunican sistemas construidos con tecnologías distintas y cómo se controla la ejecución de modelos que pueden depender de entornos complejos. Por esta razón, este capítulo presenta los fundamentos teóricos necesarios para contextualizar el problema y justificar la solución planteada.

El marco teórico se organiza en siete bloques relacionados de forma progresiva. Primero, se revisa la anotación multimodal y el papel de ELAN en la construcción de corpus audiovisuales. A partir de esa base, se analiza el reconocimiento automático como apoyo a tareas de segmentación y etiquetado. Luego, se estudian arquitecturas de integración que permiten conectar aplicaciones de escritorio con servicios externos, dando paso al análisis de servicios REST y contratos JSON como mecanismos de comunicación. Después, se revisa el uso de contenedores y la ejecución controlada de modelos, ya que estos conceptos ayudan a explicar por qué es necesario aislar dependencias y gestionar recursos. Finalmente, se aborda la orquestación de modelos de inteligencia artificial y las consideraciones de interoperabilidad, trazabilidad y mantenibilidad que sostienen la propuesta del TIC.

\subsubsection{Herramientas de anotación multimodal}

Para comenzar, la anotación multimodal consiste en asociar descripciones lingüísticas, gestuales, acústicas o visuales con intervalos específicos de tiempo dentro de una señal de audio o video. En investigaciones sobre interacción humana, gestos, conversación y lenguas de señas, la anotación no se limita a transcribir palabras, sino que también permite registrar eventos sincronizados en el tiempo, relaciones entre capas de análisis y metadatos que facilitan recuperar, comparar y reutilizar un corpus. Por esta razón, las herramientas de anotación deben ofrecer precisión temporal, soporte para múltiples niveles de descripción y capacidad de vincular anotaciones con diferentes tipos de archivos multimedia \cite{wittenburg2006elan}.

Dentro de este tipo de herramientas, ELAN tiene un papel importante porque permite trabajar con videos, audios y anotaciones organizadas por capas o \textit{tiers}. Cada capa puede representar un aspecto diferente del análisis, por ejemplo glosas, comentarios, traducciones, rasgos no manuales o marcas de revisión. Esta organización es útil para el estudio de lenguas de señas, ya que en una misma secuencia pueden ocurrir movimientos de manos, expresiones faciales, orientación corporal y otros elementos visuales que deben ser observados de manera sincronizada.

Esta organización por capas se apoya en el formato EAF, que es el formato nativo utilizado por ELAN para representar anotaciones mediante estructuras temporales. Conceptualmente, un archivo EAF permite vincular valores textuales con marcas de inicio y fin, por lo que cada anotación queda asociada a una posición específica dentro del material audiovisual. Su utilidad principal no se encuentra solo en almacenar texto, sino en conservar la relación entre el evento observado y su ubicación temporal. Este punto es importante porque cualquier apoyo automático a la anotación debe respetar esa relación para que el resultado pueda ser revisado por un especialista.

La literatura sobre ELAN destaca su flexibilidad y su énfasis en la exactitud temporal \cite{wittenburg2006elan}. Además, existen antecedentes donde se ha buscado integrar detectores de sonido e imagen dentro del flujo de anotación, como ocurrió con el proyecto AVATecH, orientado a incorporar componentes de reconocimiento de patrones en procesos de anotación semiautomática \cite{auer2010elan}. Este antecedente permite entender que la anotación asistida no pretende reemplazar el criterio del investigador, sino entregar propuestas de segmentos o etiquetas que luego deben ser evaluadas.

En el contexto de la Lengua de Señas Ecuatoriana (LSEC), esta necesidad se vuelve más evidente. La LSEC es una lengua visual-gestual utilizada por la comunidad sorda ecuatoriana y su documentación requiere conservar la relación entre movimiento, temporalidad y significado. Diferentes estudios han señalado la importancia de la LSEC en procesos de inclusión educativa y comunicación accesible \cite{ureta2022lsec_inclusion}. Asimismo, trabajos previos de la Escuela Politécnica Nacional han explorado el uso de técnicas de visión artificial y aprendizaje profundo para interpretar frases de LSEC, lo que evidencia un interés creciente por relacionar corpus visuales con modelos computacionales \cite{guevara2022lsec_interprete}.

A partir de esta relación entre ELAN, corpus audiovisuales y LSEC, la glosa se entiende como una etiqueta de trabajo que permite identificar una seña dentro de un intervalo temporal. Una glosa no representa toda la riqueza lingüística de una seña, pero permite indexar, buscar y entrenar modelos sobre segmentos previamente anotados. En consecuencia, la anotación manual sigue siendo una referencia esencial para validar cualquier resultado automático. La asistencia computacional puede reducir trabajo repetitivo o proponer intervalos iniciales, pero la decisión final sobre la corrección de una etiqueta depende del criterio humano y de la calidad del corpus disponible.

Por último, los sistemas de anotación semiautomática deben considerar que el material audiovisual suele ser heterogéneo. La variación en iluminación, encuadre, velocidad de ejecución, morfología del señante, cámara, fondo y disponibilidad de metadatos puede afectar la confiabilidad de los resultados. Por esta razón, el marco teórico no solo debe explicar cómo se almacenan las anotaciones, sino también por qué los modelos automáticos necesitan ser tratados como herramientas de apoyo. El objetivo conceptual no es convertir a ELAN en un sistema de inteligencia artificial, sino permitir que una herramienta de anotación consolidada pueda relacionarse con servicios externos especializados.

% Sugerencia de figura opcional: un esquema conceptual de ELAN con video, tiers y anotaciones temporales puede ayudar en esta subsección. Si se incluye, conviene que sea una figura simple y conceptual, no un diagrama de implementación.

\subsubsection{Reconocimiento automático en procesos de anotación}

Una vez establecido el papel de ELAN y de la anotación multimodal, es necesario revisar el concepto de reconocimiento automático. En videos de lenguas de señas, el reconocimiento automático consiste en extraer patrones de una señal visual para producir inferencias sobre eventos, acciones, gestos o categorías. En un proceso de anotación lingüística, estas inferencias pueden expresarse como segmentos temporales, etiquetas candidatas, puntuaciones de confianza o listas de alternativas. A diferencia de la anotación manual, que depende de observación experta, el reconocimiento automático se apoya en modelos estadísticos o de aprendizaje profundo entrenados con datos previos.

El primer problema asociado a este tipo de reconocimiento es la segmentación temporal. En archivos multimedia, segmentar significa identificar los intervalos donde ocurre un evento relevante. En lengua de señas, esto puede corresponder al inicio y fin de una seña, de un gesto o de una frase signada. Después de identificar el intervalo, aparece un segundo problema: la clasificación. Clasificar consiste en asignar una etiqueta al segmento detectado. Estas dos tareas se relacionan entre sí, pero no son iguales; por ello, pueden abordarse con un solo modelo o mediante una cadena de modelos donde primero se detectan intervalos y luego se clasifican.

Esta separación entre segmentación y clasificación resulta útil porque permite analizar errores de forma más clara. Un sistema puede detectar correctamente un intervalo, pero asignar una glosa incorrecta; también puede clasificar bien un fragmento, pero equivocarse al definir sus límites temporales. En corpus limitados, esta separación ayuda a tratar de forma independiente el problema temporal y el problema lingüístico. Además, permite que el especialista revise los resultados desde dos perspectivas: si el segmento está bien delimitado y si la etiqueta propuesta es adecuada.

La anotación asistida se ubica justamente entre el trabajo manual y la automatización completa. El sistema propone segmentos o etiquetas, y el especialista valida, corrige o descarta los resultados obtenidos. Este enfoque es importante para la creación de corpus audiovisuales porque los modelos pueden acelerar tareas repetitivas, pero no eliminan la necesidad de revisión. La validación humana sigue siendo necesaria por la variabilidad natural de las señas, la posible ambigüedad entre glosas, la dependencia del contexto y la presencia de rasgos no manuales que pueden no ser capturados por un modelo específico.

Para que el reconocimiento automático pueda operar sobre videos, es necesario definir cómo se representará la información visual. Un modelo puede trabajar directamente con frames del video o con características extraídas a partir de esos frames. El procesamiento basado en frames conserva información visual amplia, pero puede aumentar el costo computacional y requerir arquitecturas capaces de manejar dimensiones espaciales y temporales. En cambio, el procesamiento basado en puntos clave convierte el video en coordenadas corporales o manuales, lo que puede reducir ruido de fondo y facilitar el análisis secuencial. MediaPipe se ubica en esta segunda línea, ya que proporciona una infraestructura para construir pipelines de percepción y extraer landmarks corporales, faciales o de manos \cite{lugaresi2019mediapipe}.

Cuando la información se organiza como una secuencia, los modelos temporales adquieren mayor relevancia. Las redes LSTM fueron propuestas para modelar dependencias en secuencias y mitigar problemas asociados al aprendizaje de relaciones de largo plazo \cite{hochreiter1997lstm}. En procesamiento de señas, esta capacidad es importante porque la identidad de una seña no depende únicamente de una postura estática, sino también del movimiento realizado antes y después de cada instante. Las variantes bidireccionales permiten considerar información contextual en ambas direcciones temporales, lo que resulta útil cuando se busca segmentar eventos dentro de una secuencia completa.

Con esta base secuencial, el esquema BIO permite representar la ubicación de unidades dentro de una serie temporal. Este esquema, utilizado originalmente en tareas de segmentación de texto, codifica posiciones mediante etiquetas que indican inicio, interior o exterior de una unidad \cite{ramshaw1995text_chunking}. En video, el mismo principio puede adaptarse para marcar frames como inicio de segmento, continuación de segmento o fondo. Sin embargo, estas predicciones suelen requerir postprocesamiento para corregir secuencias inconsistentes, unir fragmentos cercanos, eliminar segmentos demasiado cortos y calcular valores de confianza.

Después de la segmentación, la clasificación de glosas puede apoyarse en modelos capaces de analizar relaciones entre elementos de una secuencia. Los modelos basados en Transformer introdujeron mecanismos de atención que permiten ponderar relaciones entre distintos elementos sin depender exclusivamente de recurrencia \cite{vaswani2017attention}. En clasificación de glosas, una arquitectura de este tipo puede procesar una secuencia de puntos clave y producir una distribución de probabilidad sobre un vocabulario. Esto no significa que el modelo comprenda todos los aspectos lingüísticos de la seña, sino que aprende patrones asociados a etiquetas observadas durante el entrenamiento.

La disponibilidad de datos útiles es una limitante crítica en lenguaje de señas. Cuando un corpus contiene pocas muestras por glosa, segmentaciones no balanceadas o anotaciones incompletas, el modelo puede aprender patrones específicos del conjunto disponible y generalizar de forma limitada. Por esta razón, las salidas automáticas deben presentarse como apoyo y no como sustituto de la anotación experta. En términos conceptuales, esto justifica que una respuesta de reconocimiento incluya no solo una etiqueta final, sino también confianza, alternativas y trazabilidad suficiente para que el especialista evalúe el resultado.

En síntesis, el reconocimiento automático aporta una base para asistir la anotación, pero no elimina la complejidad lingüística del análisis. La segmentación, la clasificación, la extracción de características y el postprocesamiento son conceptos que ayudan a comprender por qué un middleware de apoyo debe manejar resultados parciales, errores y respuestas revisables. La explicación detallada de cómo se organiza este flujo en el sistema desarrollado corresponde al capítulo de metodología.

% Sugerencia de figura opcional: se puede incluir un diagrama conceptual de segmentación temporal frente a clasificación de glosas. Esta figura ayudaría si el lector no está familiarizado con la diferencia entre detectar intervalos y asignar etiquetas.

\subsubsection{Arquitecturas de integración}

Luego de revisar la anotación y el reconocimiento automático, el siguiente problema teórico es la integración entre herramientas construidas con propósitos y tecnologías diferentes. Las aplicaciones de escritorio consolidadas suelen concentrar flujos de trabajo especializados, interfaces maduras y formatos de archivo propios. Sin embargo, su arquitectura puede no estar preparada para incorporar dependencias modernas de inteligencia artificial, especialmente cuando estas dependen de lenguajes, bibliotecas nativas, aceleración por hardware o entornos de ejecución distintos. Bajo este contexto, integrar modelos directamente dentro de la aplicación principal puede aumentar el acoplamiento, dificultar el mantenimiento y comprometer la estabilidad del software.

Una alternativa a nivel de arquitectura de software consiste en implementar una capa intermedia o middleware. En términos generales, un middleware actúa como intermediario entre sistemas que no comparten la misma interfaz, lenguaje, protocolo o ciclo de vida. Su función no se limita a reenviar mensajes, sino que también puede validar datos, traducir contratos, gestionar errores, seleccionar servicios, registrar eventos y aislar dependencias. En el presente TIC, este concepto se relaciona con la necesidad de conectar una herramienta de anotación construida sobre el ecosistema Java con servicios de inteligencia artificial que normalmente se ejecutan en entornos Python.

Dentro de esta forma de integración, el patrón \textit{sidecar} permite entender cómo un componente auxiliar puede acompañar a una aplicación principal. En arquitecturas basadas en contenedores, este patrón describe servicios complementarios que proporcionan funcionalidades de soporte sin modificar de forma invasiva la lógica central de la aplicación \cite{burns2016container_patterns}. En el contexto de una aplicación de escritorio, el principio se mantiene porque el componente auxiliar ejecuta tareas especializadas y expone una interfaz controlada, mientras que la aplicación principal conserva su flujo de trabajo y delega únicamente las operaciones externas necesarias.

Esta separación ayuda a que cada componente mantenga una responsabilidad clara. ELAN conserva el rol de herramienta de anotación y visualización, mientras que el middleware se concibe como una capa de orquestación local encargada de recibir solicitudes, validar parámetros, seleccionar servicios y devolver resultados estructurados. De esta manera, los modelos de inteligencia artificial pueden evolucionar sin exigir que la aplicación principal incorpore de forma nativa todas sus dependencias. Además, se reduce el riesgo de que una falla de inferencia o una incompatibilidad de bibliotecas afecte directamente el proceso de anotación.

La integración entre aplicaciones de escritorio y servicios externos puede realizarse mediante archivos intermedios, llamadas de sistema, sockets, servicios HTTP u otros mecanismos de comunicación. La elección depende de requisitos como portabilidad, facilidad de depuración, seguridad, rendimiento y soporte de bibliotecas. Para un entorno local, un servicio HTTP sobre la dirección local ofrece una frontera clara entre procesos y permite reutilizar herramientas comunes de prueba, documentación y validación. Esta decisión también favorece la independencia entre lenguajes, dado que la aplicación Java y el servicio Python solo necesitan seguir el contrato de comunicación establecido.

El desacoplamiento es un principio central en este tipo de arquitectura. Cuando el cliente, el middleware y los modelos comparten contratos explícitos, cada parte puede evolucionar sin impactar de forma directa sobre las demás. Desde la perspectiva del cliente, no debería ser necesario conocer si la inferencia se ejecuta mediante una simulación, una ejecución nativa o un entorno aislado. Del mismo modo, el middleware puede incorporar nuevos modelos siempre que estos respeten los contratos de entrada y salida previamente definidos para el sistema.

La arquitectura de integración también debe considerar la experiencia del usuario final. En un flujo de anotación, el investigador espera seleccionar material audiovisual, activar una tarea de reconocimiento y recibir segmentos editables dentro de su entorno de trabajo. Por esta razón, el middleware no debe imponer un cambio conceptual sobre la tarea lingüística. Su valor consiste en ocultar complejidad técnica y devolver resultados compatibles con las estructuras de anotación utilizadas por ELAN. De esta manera, la innovación se concentra en la interoperabilidad y no en reemplazar la herramienta original.

\subsubsection{Servicios REST y contratos de comunicación}

La necesidad de integración descrita en la sección anterior conduce al análisis de los servicios REST y los contratos de comunicación. REST, propuesto por Fielding como estilo arquitectónico para sistemas distribuidos, define principios para organizar recursos, representaciones y operaciones sobre una interfaz uniforme \cite{fielding2000rest}. En un servicio REST, las funcionalidades se exponen mediante endpoints, pueden utilizarse con métodos HTTP y las respuestas se comunican mediante códigos de estado y representaciones estructuradas. Aunque REST no establece un formato de datos único, JSON se utiliza ampliamente por su compatibilidad con múltiples lenguajes y por su facilidad para representar objetos, listas, cadenas, números y valores booleanos.

En una integración entre ELAN y un middleware de inteligencia artificial, una API REST permite convertir una solicitud de anotación en un mensaje explícito. El cliente puede enviar la ruta del video, la configuración del modelo, el tier destino, parámetros de segmentación y preferencias de ejecución. A su vez, el servicio puede responder con estado del proceso, información del medio, segmentos temporales y trazas de ejecución. Esta forma de comunicación evita acoplar directamente objetos internos de Java con objetos internos de Python, y también permite probar el servicio mediante clientes HTTP independientes antes de integrarlo con la aplicación de escritorio final.

Para que esta comunicación sea confiable, los contratos deben estar definidos con claridad. Un contrato de comunicación especifica campos obligatorios, tipos de datos, rangos permitidos, estructura de respuesta y errores esperados. En un sistema de anotación asistida, los campos \texttt{start\_ms}, \texttt{end\_ms}, \texttt{label} y \texttt{confidence} representan el núcleo de un segmento anotable. Los dos primeros delimitan el intervalo temporal, el tercero contiene la etiqueta propuesta y el cuarto expresa una medida de confianza. Esta estructura permite que un cliente convierta resultados automáticos en anotaciones temporales visibles en la línea de tiempo donde trabaja el lingüista.

La validación de esquemas complementa esta idea de contrato. FastAPI proporciona una infraestructura para declarar endpoints, asociar modelos de datos y generar documentación OpenAPI a partir de tipos de Python \cite{fastapi2026features}. Pydantic, por su parte, permite definir modelos de validación y conversión de datos, de modo que un payload incorrecto pueda rechazarse antes de ejecutar lógica de negocio \cite{pydantic2026models}. En un middleware de inferencia, esta validación es importante porque evita ejecutar procesos costosos cuando faltan campos, los tipos no coinciden o los parámetros no cumplen restricciones mínimas.

Los códigos de estado y las respuestas de error también forman parte del contrato. Una API robusta no debe limitarse a devolver datos cuando todo funciona, sino que también debe informar si el video no existe, si el modelo no está registrado, si una arquitectura no coincide, si el procesamiento con CUDA fue solicitado pero no está disponible o si el vocabulario de glosas no cumple el formato esperado. Las respuestas estructuradas facilitan la depuración y permiten que el cliente presente mensajes controlados sin cerrar la aplicación principal. Esto resulta especialmente útil cuando el servicio forma parte de un flujo de anotación interactivo.

En este punto, AVATecH aporta un antecedente relevante para la comunicación entre ELAN y reconocedores externos. Su especificación técnica describe mecanismos para que extensiones de reconocimiento intercambien parámetros, archivos y resultados con ELAN \cite{auer2014avatech}. Aunque la tecnología y los modelos actuales han evolucionado, el principio se mantiene: la herramienta de anotación necesita una frontera bien definida para invocar procesos externos y recibir anotaciones candidatas. Por esta razón, la traducción conceptual de un protocolo heredado hacia una API REST moderna debe preservar la semántica temporal y la estructura de los resultados.

Además de los datos de entrada y salida, los metadatos influyen en la interoperabilidad. CMDI fue propuesto como un marco basado en componentes y categorías de datos para mejorar la consistencia de metadatos en recursos lingüísticos \cite{broeder2010cmdi}. En un sistema de anotación asistida, los metadatos permiten describir modelos, entradas, salidas, parámetros y recursos asociados. Esto evita que los resultados generados por un componente automático queden aislados, ya que pueden relacionarse con el corpus, la herramienta de anotación y la configuración de ejecución.

\subsubsection{Contenedores y ejecución controlada de modelos}

Después de definir la comunicación entre sistemas, es necesario revisar cómo se controla el entorno donde se ejecutan los modelos. La contenedorización permite empaquetar una aplicación junto con dependencias, bibliotecas y configuración de ejecución en una unidad portable. Docker describe los contenedores como entornos aislados que ejecutan procesos sobre el sistema anfitrión y que pueden construirse a partir de imágenes reproducibles \cite{docker2026overview}. Para sistemas de inteligencia artificial, este enfoque resulta atractivo porque los modelos suelen depender de versiones específicas de Python, bibliotecas numéricas, frameworks de aprendizaje profundo, controladores y recursos de hardware.

En una arquitectura que integra ELAN con modelos de inteligencia artificial, el uso de contenedores ofrece una separación clara entre la herramienta de anotación y el entorno de inferencia. ELAN no necesita instalar directamente PyTorch, OpenCV, MediaPipe u otras dependencias nativas. El modelo puede ejecutarse en un entorno controlado y exponer resultados mediante el middleware. Esta separación reduce conflictos de versiones y permite reproducir el comportamiento del servicio en distintos equipos, siempre que se mantenga la misma imagen o configuración de ejecución.

La ejecución nativa y la ejecución contenerizada representan dos estrategias distintas. En la ejecución nativa, el servicio carga las bibliotecas directamente desde el entorno local. Esta modalidad suele ser útil durante el desarrollo porque permite depurar con menor sobrecarga. En la ejecución contenerizada, el modelo se ejecuta en un entorno aislado, lo que mejora la portabilidad y el control de dependencias, pero también introduce aspectos adicionales como montaje de archivos, comunicación con el contenedor, acceso a GPU, límites de recursos y ciclo de vida del proceso.

El uso de contenedores no elimina la necesidad de contratos. Un modelo contenerizado debe recibir entradas en un formato acordado y devolver salidas compatibles con el middleware. Por ello, la contenedorización debe entenderse como una estrategia de aislamiento y reproducibilidad, no como un sustituto de una arquitectura de comunicación. El middleware sigue siendo responsable de validar solicitudes, seleccionar el mecanismo de ejecución adecuado, controlar errores y convertir resultados a una representación que ELAN pueda utilizar.

La ejecución controlada de modelos también se relaciona con recursos de cómputo. Los modelos de video pueden requerir memoria considerable, aceleración por GPU o tiempos de inferencia altos. Un sistema de orquestación debe contemplar parámetros como dispositivo preferido, tiempo máximo de ejecución y manejo de fallos. Las métricas de consumo, latencia y rendimiento son fundamentales para evaluar el comportamiento del sistema, pero sus valores concretos corresponden al capítulo de resultados porque dependen de pruebas experimentales y no del marco teórico.

Desde el punto de vista académico y técnico, el uso de contenedores mejora la reproducibilidad en diferentes entornos y ayuda a separar responsabilidades. La herramienta de anotación conserva su entorno nativo, el middleware coordina solicitudes y los modelos se ejecutan bajo condiciones definidas. Esta organización es coherente con el principio de mantener bajo acoplamiento entre componentes y de introducir cambios técnicos sin alterar innecesariamente la experiencia del usuario final.

% Sugerencia de figura opcional: si se desea agregar una imagen en esta subsección, conviene que sea un diagrama conceptual sobre aislamiento de dependencias. La arquitectura detallada de contenedores debería reservarse para metodología.

\subsubsection{Orquestación de modelos de inteligencia artificial}

Una vez revisada la ejecución controlada, el concepto que articula el flujo completo es la orquestación de modelos de inteligencia artificial. Orquestar no significa únicamente llamar a un modelo, sino coordinar las etapas necesarias para ejecutar una inferencia de forma controlada. Esto incluye registrar modelos disponibles, validar artefactos, preparar entradas, seleccionar el mecanismo de ejecución, controlar dispositivo, ejecutar inferencia, transformar salidas, postprocesar resultados y registrar trazas. En un contexto de anotación, la orquestación también debe asegurar que la salida final sea interpretable como segmentos temporales o etiquetas candidatas.

Un registro de modelos es un componente conceptual que almacena información sobre modelos disponibles, versiones, estado, tipo de tarea, artefactos y contratos de entrada y salida. Su utilidad radica en separar la identidad del modelo de su ejecución concreta. El cliente puede solicitar un \texttt{model\_id} y una versión, mientras que el middleware resuelve dónde se encuentran los artefactos, qué mecanismo de ejecución debe utilizarse y qué restricciones aplican. Este patrón favorece la extensibilidad porque nuevos modelos pueden incorporarse sin cambiar la forma general de la solicitud.

La selección del mecanismo de ejecución es otra función central dentro de la orquestación. Un runner o ejecutor encapsula la forma específica de ejecutar un modelo, una simulación o un pipeline compuesto. Esta abstracción evita que la capa que recibe solicitudes tenga que conocer los detalles de cada arquitectura. En términos de diseño, el ejecutor traduce una entrada interna normalizada a la forma que necesita el modelo y devuelve una salida que luego puede ser transformada por el postprocesamiento.

PyTorch es un framework ampliamente utilizado para aprendizaje profundo, diseñado con un estilo imperativo y soporte para cómputo tensorial acelerado \cite{paszke2019pytorch}. En este trabajo de integración, su importancia conceptual está ligada con la ejecución de modelos de video y secuencia dentro de un entorno Python. Sin embargo, la existencia de PyTorch como dependencia no implica que cualquier arquitectura sea compatible automáticamente. Se requieren pesos, estructuras de modelo y formas de tensores coherentes con los contratos definidos por el middleware.

Antes de la inferencia, el preprocesamiento convierte el video original en una representación adecuada para el modelo. Puede incluir lectura de metadatos, muestreo de frames, conversión de espacios de color, redimensionamiento, normalización y construcción de ventanas temporales. Estas etapas son necesarias porque los modelos de aprendizaje profundo suelen requerir tensores de forma fija y rangos numéricos específicos. En el caso de modelos basados en puntos clave, el preprocesamiento también puede incluir extracción de landmarks, normalización espacial y cálculo de características dinámicas.

Después de la inferencia, la salida del modelo debe transformarse en información útil para la anotación. Un segmentador puede devolver probabilidades por frame; un clasificador puede devolver logits o probabilidades por clase; y un pipeline compuesto puede devolver listas de alternativas. Por esta razón, el postprocesamiento es indispensable. En segmentación temporal, el postprocesamiento convierte probabilidades o etiquetas en intervalos con inicio, fin y confianza. En clasificación, selecciona una etiqueta principal y puede conservar alternativas para revisión.

El uso combinado de segmentación BIO y clasificación de glosas responde a una división conceptual entre detección temporal y asignación léxica. La etapa BIO identifica regiones candidatas en una secuencia, mientras que la etapa de clasificación asigna una glosa a cada región. Esta organización permite producir segmentos cuyo campo \texttt{label} contenga una glosa propuesta y cuyo campo \texttt{predictions} incluya alternativas ordenadas. La presencia de alternativas no debe interpretarse como certeza lingüística, sino como información adicional para la revisión humana.

La trazabilidad de ejecución permite reconstruir lo que ocurre durante una solicitud. En un sistema de orquestación, una traza puede incluir mecanismo de ejecución seleccionado, dispositivo usado, modelo, versión, número de ventanas, estados atravesados y tiempos por etapa. Estas marcas no sustituyen una evaluación experimental, pero facilitan depuración, auditoría técnica y comparación controlada entre ejecuciones. Cuando un modelo falla, la traza y el código de error ayudan a distinguir si el problema está en la carga del modelo, lectura del video, extracción de puntos clave, inferencia o postprocesamiento.

El diseño de la orquestación debe mantener una frontera clara entre salida interna y salida externa. Internamente, un ejecutor puede trabajar con probabilidades por frame, tensores, logits o listas de predicciones. Externamente, el cliente requiere una respuesta estable y orientada a anotación. Esta diferencia evita que ELAN dependa de detalles específicos de PyTorch, MediaPipe o una arquitectura particular. El middleware opera como traductor entre el mundo de los modelos y el mundo de las anotaciones temporales.

\subsubsection{Interoperabilidad, trazabilidad y mantenibilidad}

Finalmente, los conceptos anteriores se conectan mediante tres criterios transversales: interoperabilidad, trazabilidad y mantenibilidad. La interoperabilidad se refiere a la capacidad de sistemas distintos para intercambiar información y utilizarla de forma significativa. En este trabajo de integración, la interoperabilidad no consiste únicamente en enviar datos desde ELAN hacia un servicio HTTP, sino también en conservar el significado temporal de los segmentos, mantener etiquetas compatibles con tiers de anotación y devolver errores comprensibles para el flujo de trabajo. Para lograrlo, se requieren contratos claros, documentación técnica y separación de responsabilidades entre cliente, middleware y modelos.

La compatibilidad con ELAN exige respetar la naturaleza temporal de la anotación. Un resultado automático solo es útil si puede ubicarse en la línea de tiempo y revisarse junto con el video. Por ello, los segmentos generados por un sistema de apoyo deben expresar intervalos en milisegundos y etiquetas textuales. En este marco conceptual, la generación directa de archivos EAF no es el punto central; lo importante es que la información producida mantenga la semántica temporal necesaria para convertirse en anotaciones revisables dentro de ELAN.

La trazabilidad técnica se apoya en logs, estados, métricas internas y respuestas estructuradas. En un sistema de inferencia, una falla puede originarse por múltiples causas, como un archivo de video inexistente, un paquete de modelo inválido, un vocabulario incompleto, incompatibilidad de pesos, ausencia de procesamiento CUDA, salida con forma inesperada o error de extracción de puntos clave. Sin trazabilidad, estas fallas se vuelven difíciles de diagnosticar. Con trazabilidad, el sistema puede informar un código específico y orientar la corrección sin interrumpir la herramienta principal.

La mantenibilidad depende de que los módulos tengan responsabilidades delimitadas. Separar rutas de comunicación, esquemas, servicios, almacenamiento, procesamiento de video y mecanismos de ejecución permite evolucionar una parte sin reescribir todo el sistema. Este principio es especialmente relevante cuando el proyecto necesita incorporar nuevos modelos, cambiar el mecanismo de ejecución o adaptar la comunicación con ELAN. Una arquitectura modular también facilita documentar evidencias, revisar pruebas y justificar decisiones técnicas.

La documentación técnica cumple una función de puente entre implementación y reproducibilidad académica. No basta con que el sistema funcione localmente, también debe poder explicarse, verificarse y mantenerse. Los documentos de fases, contratos de endpoints, ejemplos de solicitudes y descripciones de errores permiten que otro lector comprenda el alcance real del componente. Además, ayudan a diferenciar entre funcionalidades implementadas, funcionalidades planificadas y limitaciones pendientes.

La soberanía de los datos es una consideración adicional en corpus audiovisuales. El procesamiento local evita enviar videos sensibles o material de investigación a servicios externos no controlados. Este aspecto resulta relevante cuando se trabaja con datos de personas, señas, contextos educativos o corpus en construcción. La arquitectura de middleware local favorece que la inferencia se ejecute en el equipo del usuario o en un entorno controlado, aunque la protección efectiva también depende de políticas de almacenamiento, permisos y manejo de archivos que deben documentarse en la metodología y validación.

Por otro lado, la mantenibilidad también exige evitar dependencias implícitas. Si ELAN conoce detalles internos de un modelo, cada cambio de arquitectura obligaría a modificar la herramienta de anotación. Si el modelo asume estructuras internas de ELAN, se reduce la posibilidad de reutilizarlo en otros clientes. El middleware resuelve este problema al establecer una interfaz estable. Mientras el contrato externo se mantenga, los cambios internos pueden limitarse al registro de modelos, mecanismos de ejecución o postprocesadores.

En conclusión, el marco teórico permite entender que la solución propuesta no se sustenta en un único concepto aislado. La anotación multimodal explica la necesidad de conservar intervalos temporales; el reconocimiento automático justifica el uso de modelos para proponer segmentos y etiquetas; las arquitecturas de integración permiten conectar ELAN con servicios externos; REST y JSON ofrecen contratos de comunicación; la contenedorización aporta aislamiento; y la orquestación permite coordinar modelos, entradas, salidas y trazas. Estos fundamentos establecen la base conceptual para que la metodología explique, en el siguiente capítulo, cómo se materializó la solución dentro del TIC.

